(1) 取 Fourier 表示 \(A(z,t)=\int\widetilde A(z,\Omega)e^{-i\Omega t}d\Omega/(2\pi)\)。因 \(\partial_t^2\mapsto-\Omega^2\)，方程化为
\[
\partial_z\widetilde A=i\beta_2\Omega^2\widetilde A/2,\quad
\widetilde A(z,\Omega)=\widetilde A(0,\Omega)e^{i\beta_2z\Omega^2/2}.
\]
谱幅度模长不变，仅积累二次色散相位。

(2) 初始 Fourier 高斯 \(\widetilde A(0,\Omega)=A_0T_0\sqrt{2\pi}e^{-T_0^2\Omega^2/2}\)。令 \(s=\beta_2z/T_0^2\)，完成 Gaussian 积分得到
\[
A(z,t)=\frac{A_0}{\sqrt{1-is}}
\exp\left[-\frac{t^2}{2T_0^2(1-is)}\right].
\]
平方根分支从 z=0 的正值连续选取。

(3) 输出强度
\[
|A(z,t)|^2=\frac{|A_0|^2}{\sqrt{1+s^2}}
\exp[-t^2/(T_0^2(1+s^2))].
\]
所以强度 \(1/e\) 半宽 \(T(z)=T_0\sqrt{1+s^2}\)。积分全时间，峰值前因子 \(1/\sqrt{1+s^2}\) 与宽度 \(\sqrt{1+s^2}\) 抵消，故 \(\int|A|^2dt=|A_0|^2T_0\sqrt\pi\) 守恒。

(4) \(\beta_2=0\) 时 s=0，脉冲不展宽；若 \(|\beta_2z|\gg T_0^2\)，则 \(T(z)\sim|\beta_2|z/T_0\)。色散展宽是频率分量群时延不同的结果，不需要吸收；本结论仅适用于二阶色散和线性传播模型。
